\section{Fixed multiples of the girth and bounded-rank exceptions}
\label{sec:four-high-connectivity-weighted-kernel}

Suppressing the ports gives a cubic kernel with positive integer edge lengths.
A weighted Moore argument isolates the possible bounded-rank exceptions.
The path-stability agent proposed the walk argument below.
All logarithms have base two.

\begin{lemma}[A weighted-kernel Moore bound]
\label{lem:four-weighted-kernel-moore}
Let $N$ be a normalized $k$-pole of rank $h\ge2$ and girth $g\ge3$.
Then, with $t=\lceil\log(8h)\rceil$,
\[
 g\le4t\left(1+\frac{k}{3(h-1)}\right).
\]
\end{lemma}

\begin{proof}
Suppress all degree-two vertices.
Since $N$ is connected and has rank at least two, the resulting kernel is a connected cubic multigraph.
Loops and parallel edges are allowed.
It has $n_0=2h-2$ vertices and $m_0=3h-3$ edges.
Give each kernel edge the length of the path that it replaces.
These positive lengths sum to $m_0+k$.
The weighted girth of the kernel equals $g$.

There are $2m_0\,2^{t-1}$ directed non-backtracking walks of length $t$ in the kernel.
Choose one uniformly.
Every directed edge is uniform at every step, so its expected total weight is
\[
 t\left(1+\frac{k}{m_0}\right).
\]
At least half the walks have weight at most twice this expectation.
Their number is at least
\[
 m_0\,2^{t-1}\ge12h(h-1)>(2h-2)^2=n_0^2.
\]

If one such walk repeats a vertex, its edge support contains a cycle of no greater weight.
Otherwise these walks are simple paths.
Two must have the same ordered endpoints, so their union contains a cycle of weight at most the sum of their weights.
In either case, a cycle has weight at most $4t(1+k/m_0)$.
This proves the bound.
\end{proof}

\begin{proposition}[A bounded-rank or near-Moore alternative]
\label{prop:four-fixed-ratio-shore-rank}
Fix $C>0$, and put
\[
 H_C=\left\lceil64(C+1)\log(16(C+1))\right\rceil+1.
\]
If a normalized $k$-pole has rank $h\ge H_C$, girth at least $g$, and $k\le Cg$, then
\[
 h\ge2^{g/8-4}.
\]
For sufficiently large $g$, depending only on $C$, the stronger bound
\[
 h\ge2^{g/2-5}
\]
holds.
\end{proposition}

\begin{proof}
The choice of $H_C$ ensures
\[
 h-1>\frac{8C}{3}\,\lceil\log(8h)\rceil
 \qquad(h\ge H_C).
\]
Here is one explicit verification.
Put $D=C+1$ and $u=\log(16D)\ge4$.
At $h=H_C$, one has $h-1\ge64Du$ and $\log h+4\le4u$.
The function $(h-1)/(\log h+4)$ increases for $h\ge2$.
Also $\lceil\log(8h)\rceil\le\log h+4$.
These facts give the displayed inequality for all $h\ge H_C$.

Lemma~\ref{lem:four-weighted-kernel-moore} and $k\le Cg$ now give
\[
 g\le4t+\frac{4tCg}{3(h-1)}<4t+\frac g2.
\]
Thus $g<8t\le8\log h+32$, proving the first bound.

For sufficiently large $g$, that bound implies $h-1\ge(C+1)g^2$.
If $h\ge2^{g/2}$, the stronger assertion already holds.
Otherwise $n=2h+k-2\le2^{g/2+2}$ for large $g$, and
\[
 \frac{2k\log n}{3(h-1)}
 \le\frac{2Cg(g/2+2)}{3(C+1)g^2}<1.
\]
The entropy Moore bound therefore gives $g<5+2\log n$.
Since $n>2^{(g-5)/2}$ eventually exceeds $2Cg$, one has
\[
 h=\frac{n-k+2}{2}>\frac n4>2^{g/2-5}.
\]
\end{proof}

\begin{corollary}[Uniform terminal forests above the bounded-rank range]
\label{cor:four-fixed-ratio-terminal-forests}
For each fixed $C>0$, there is $B_C$ such that every normalized $k$-pole of rank $h\ge H_C$ with
\[
 2\le k\le Cg,
 \qquad \operatorname{girth}(N)\ge g
\]
satisfies, for every even port subset $T$,
\[
 \frac{F_T(N)}{2^h}
 \le\frac12+B_C g^{13/2}2^{-g/4}.
\]
\end{corollary}

\begin{proof}
Use Proposition~\ref{prop:four-fixed-ratio-shore-rank} in the proof of Theorem~\ref{thm:path-high-girth-terminal-forests}.
The port-free ball with radius $\lfloor g/2-\log k-10\rfloor$ still exists.
Its layer lower bound is at least $2^{g/2}/(2^{11}Cg^2)$.
The thin-marker layer upper bound has polynomial term at most $(Cg+1)^2$.
The same inversion gives the result after changing its constant by a factor depending only on $C$.
Increasing that constant handles the remaining finite range of $g$.
\end{proof}

\begin{corollary}[All fixed multiples of the girth]
\label{cor:four-fixed-ratio-cyclic-cuts}
Fix $C>0$.
Let connected cubic graphs have girths $g\to\infty$.
If a cut of order at most $Cg$ has two connected shores, each of rank at least $H_C$, then
\[
 c(G)\le\left(\frac14+o_C(1)\right)2^{\beta(G)}.
\]
\end{corollary}

\begin{proof}
Move vertices with two cut incidences across the cut until its endpoints are distinct on each shore.
Such a vertex is a leaf in its connected positive-rank shore; moving it preserves that rank, increases the other rank by one, and decreases the cut order by one.
A vertex with three cut incidences would be isolated in its shore, which is impossible.
The resulting connected normalized poles have ranks $h,j\ge H_C$ and cut order $k'\le Cg$.
Proposition~\ref{prop:four-fixed-ratio-shore-rank} makes both ranks at least $2^{g/2-5}$ for large $g$.
If $k'=1$, there are no crossing cycles and the conclusion follows from the internal-cycle bound.
Otherwise apply Corollary~\ref{cor:four-fixed-ratio-terminal-forests} to both shores and sum over the $2^{k'-1}-1$ nonempty even boundary subsets.
The crossing count is at most $(1/2+o_C(1))^2 2^{\beta(G)}$.
The internal counts, at most $2^h+2^j$, are negligible after normalization.
\end{proof}

Consequently, an unresolved above-half family with girth tending to infinity cannot have such a cut.
For every fixed $C$, each cut of order at most $Cg$ with connected shores must have a shore of rank below $H_C$.
This does not prove that an appropriate moderate cut exists in every remaining graph.
